\scanpage{63}
\mspage{27}
\deleted{c) Si $(z_\alpha)_{\alpha\in I}$ est l'ens des éléments de
$\Sigma_{<x}=\uline{b}^{-1}(\{x\})$ (où $\uline{b}:\Sigma^{\unsure{<|*}}(1)\longrightarrow\Sigma$), les
\unsure{les};}

a) Pour \unsure{tt|$\forall$} $x\in\Sigma$, l'ens $\Sigma_{<x}$ \deleted{(fini) \unsure{est}}\inserted{est tot${}^t$}
\deleted{$\simeq\uline{b}^{-1}(\{x\}$} ordonné

b) Si $\Sigma^{\unsure{*|<}}(1)$ est le graphe de la relation
$x<y$, et $\uline{s},\uline{b}:\Sigma^{<}(1)\unsure{\rightrightarrows|=}\Sigma$ sont les applications
source et but, \deleted{\illegible{1}} l'application $\uline{b}$ est une immer-
sion \deleted{\uline{\unsure{bordante}}} \uline{propre} (donc les $\Sigma_{<x}\unsure{\ne|=}\uline{b}^{-1}(\{x\})$
sont finis)

c) Si $(z_\alpha)_{\alpha\in I}$ sont les éléments de $\Sigma_{<x}=b^{-1}(\{x\})$,
alors \inserted{la famille} \unsure{des|les} \unsure{g}ermes \inserted{en $x$} de sous-espaces \deleted{\unsure{bordants}}\inserted{de $\Sigma$}
images par $\uline{b}$ des \inserted{germes} $(\Sigma^{<}(1),z_\alpha)$, est \inserted{une famille}
\deleted{\unsure{transversaux}}\inserted{transversale}
de sous-espaces \uline{bordants}.

d) $\uline{\unsure{b|s|p_1}}:\Sigma^{<}(1)\longrightarrow\Sigma^*$ est \uline{transversal} p.r. à l'immersion
bordante $\uline{b}:\Sigma^{<}(1)\mathrel{-}\Sigma^{\unsure{*}}$
\deleted{\unsure{la diag.}}
\margin{ou sous forme d'une \unsure{récollection|récollement}
pour les $z_\alpha\in\Sigma^{<}(1)$ au dessus d'un même
$x\in\Sigma^*$\ldots}

$1^{\circ})$ Cas équisingulier : $\uline{s}:\Sigma^{<}(1)\mathrel{-}\Sigma$ fibrant \inserted{(propre)}

$2^{\circ})$ Cas non singulier : $\unsure{\Sigma}$ est une variété à bord,
dont le bord est $\uline{b}(\Sigma^{*}(1))$:

$3^{\circ})$ Cas \unsure{équisingulier|équi-singulier}-non-singulier : $\uline{s}:\Sigma^{<}(1)\mathrel{-}\Sigma$ est
fibrant \inserted{propre}, à fibres des variétés à bord, et le
fibré des bords est l'image de $\Sigma^{*}(2)$ par
le morphisme $\Sigma^{*}(2)\xrightarrow{\uline{s'}}\Sigma^{\unsure{*|<}}(1)$ \deleted{\unsure{avec}}\inserted{donné}
par $(x\prec y\prec z)\longrightarrow(x\prec z)$.
\margin{c'est sans doute conséquence de la conjonction des deux conditions $1^{\circ}),2^{\circ})$}
\margin{62}
